% =====================================================================
%  Hiện thực UC01, UC01, UC03 (chèn vào hien-thuc-uc.tex trước mục UC06)
% =====================================================================
\subsection{Đăng ký, xác thực OTP và khôi phục mật khẩu (UC01, UC01, UC03)}\label{sssec:ht-dang-ky}

Ba use case do \texttt{identity.Service} hiện thực qua năm API công khai: A01 \texttt{POST} \path{/auth/register}, A02 \texttt{POST} \path{/auth/otp/verify}, A03 \texttt{POST} \path{/auth/otp/resend}, A07 \texttt{POST} \path{/auth/password/forgot} và A08 \texttt{POST} \path{/auth/password/reset}. Người gọi chưa có danh tính để làm chủ khóa lũy đẳng, nên các API này không dùng \texttt{Idempotency-Key} mà dựa vào giới hạn tần suất, bộ đếm của \br{BR02} và tính dùng một lần của mã. Hình~\ref{fig:seq-dang-ky-otp} thể hiện luồng đăng ký và xác thực; A02 rẽ nhánh theo mục đích của mã.

\diagram[\textwidth]{c3-35-seq-dang-ky-otp}{Biểu đồ tuần tự hiện thực use case Đăng ký tài khoản, Xác thực mã OTP và Khôi phục mật khẩu}{fig:seq-dang-ky-otp}

A01 chịu giới hạn 5 lần mỗi giờ theo IP, rồi được kiểm tra \br{BR01} trước mọi truy vấn: thiếu trường trả 400 mã 2001 (\msg{MSG01}), sai định dạng trả 400 mã 2002 (\msg{MSG02}). Phép băm Argon2id tốn vài chục mili giây CPU nên chạy trước khi mở giao dịch. Bộ đếm \br{BR02} ở Redis gồm \texttt{SET NX EX 60} cho thời gian chờ gửi lại và \texttt{INCR} theo giờ cho tối đa 5 lần gửi; vi phạm trả 429 mã 4290 kèm \texttt{Retry-After} và \msg{MSG30} (bước 4--5). Ba trường hợp của email là chưa có, đang chờ xác thực và đã kích hoạt được phân biệt bằng một câu lệnh (bước 6): \texttt{INSERT ... ON CONFLICT (email) DO UPDATE ... WHERE users.status = 'pending' RETURNING id}. Không có dòng trả về nghĩa là email đã thuộc tài khoản hoạt động, dịch vụ trả 409 mã 4003 kèm \msg{MSG03}; vì chèn và cập nhật là một thao tác nguyên tử trên \texttt{uq\_users\_email}, hai yêu cầu đồng thời không thể tạo hai tài khoản chờ. \msg{MSG03} để lộ việc email đã đăng ký; đó là chủ ý của \uc{UC01} để hướng người dùng sang đăng nhập hoặc khôi phục, và mức lộ bị giới hạn theo IP chặn trên. Email đang chờ được ghi đè họ tên, mật khẩu và tiếp tục xác minh. Cùng giao dịch vô hiệu các mã chưa dùng của cặp email, mục đích rồi chèn mã mới gồm sáu chữ số từ \texttt{crypto/rand}, lưu HMAC-SHA256 với khóa bí mật ngoài CSDL, sống \texttt{otp.ttlMinutes}. Thư do tác vụ \texttt{SendEmail} gửi, được chèn bằng \texttt{InsertTx} (bước 8): SMTP chậm không kéo dài giao dịch, lỗi gửi được River thử lại~\cite{river}, còn giao dịch bị hủy thì thư không bao giờ đi. A03 lặp lại bước 4--8 cho tài khoản đang chờ và trả 200 cả khi không có tài khoản chờ, để không thành công cụ dò email.

A02 khóa dòng \texttt{otp\_challenges} mới nhất của cặp email, mục đích bằng \texttt{FOR UPDATE}, so HMAC trong thời gian hằng, kiểm tra hạn và số lần sai (bước 12); hai lần gửi cùng mã, như tự điền của hệ điều hành cộng thao tác tay, được tuần tự hóa và lần sau nhận 400 mã 3006. Nhánh mã sai là ngoại lệ có chủ đích của quy tắc “lỗi thì hủy giao dịch”: dịch vụ tăng \texttt{failed\_attempts}, đủ 5 lần thì vô hiệu mã, rồi \texttt{COMMIT} trước khi trả 400 mã 3006 kèm \msg{MSG04} (bước 13--14), vì hủy giao dịch sẽ cho kẻ dò số lượt thử không giới hạn. Với 5 lượt mỗi mã và 5 mã mỗi giờ, một email chịu tối đa 25 lượt đoán trên $10^6$ giá trị mỗi giờ, xác suất trúng khoảng $2{,}5 \cdot 10^{-5}$. Mã đúng với mục đích \texttt{register} chuyển tài khoản sang \texttt{active}, rồi \texttt{identity/api.ActivationHook} do household hiện thực chèn không gian \texttt{personal} cùng dòng thành viên vai trò \texttt{leader} trong cùng \texttt{pgx.Tx} (bước 15--17). Móc giữ đồ thị phụ thuộc không chu trình mà vẫn bảo đảm nguyên tử vế đầu của \br{BR04}: không tồn tại tài khoản hoạt động thiếu không gian cá nhân. Giao dịch đặt \texttt{last\_space\_id} và tạo phiên như A04 (bước 18), nên người dùng vào thẳng ứng dụng; nếu phản hồi bị mất, lần gửi lại nhận 3006 và ứng dụng chuyển sang đăng nhập, nơi A04 thành công. Ngược lại, tài khoản còn chờ đăng nhập đúng mật khẩu nhận 403 mã 3009 kèm \msg{MSG34}, máy chủ gửi lại mã trong giới hạn \br{BR02} để ứng dụng mở \uc{UC01}.

\uc{UC03} dùng lại cơ chế này với mục đích \texttt{reset\_password}. A07 luôn trả 200 cùng một thông điệp (Bảng~\ref{tab:UC03}, bước 2) và chỉ chèn mã cùng tác vụ \texttt{SendEmail} khi email thuộc tài khoản \texttt{active} hoặc đang khóa tạm; bộ đếm \br{BR02} tăng cho mọi email nên cả phản hồi 429 cũng không lộ email có thật. A02 với mục đích này trả \texttt{resetToken} 256 bit thay vì cấp phiên; chỉ băm SHA-256 được lưu vào \texttt{reset\_token\_hash} của dòng mã vừa dùng, sống 10 phút (bước 20--21), nhờ đó màn hình mật khẩu mới chỉ mở sau khi email đã được xác minh. A08 kiểm tra \br{BR01}, khóa dòng mã theo email và băm của token, rồi trong một giao dịch đặt mật khẩu mới, xóa \texttt{reset\_token\_hash} để token chỉ dùng một lần, đặt lại \texttt{failed\_login\_count} và \texttt{locked\_until}, thu hồi mọi phiên, tăng \texttt{token\_version}; sau commit, \path{auth:ver:<userId>} được ghi lại để mã truy cập cũ mất hiệu lực ngay (\br{BR03}). Token sai, hết hạn hoặc không khớp email trả 400 mã 3006.
